\chapter{Antecedentes}
\label{ch:antecedentes}

\pendiente{INVESTIGACIÓN PENDIENTE: este capítulo exige búsqueda de fuentes
reales y citas verificables. No rellenar con texto genérico.}

\section{Estado del arte de las plataformas de datos deportivos y los wearables}
\label{sec:antecedentes-plataformas}

\pendiente{describir el estado actual de las plataformas de datos deportivos
(Strava, TrainingPeaks, Garmin Connect y similares) y del ecosistema de
dispositivos wearables: qué métricas registran, con qué precisión y bajo qué
modelo de acceso a los datos. Cada afirmación debe apoyarse en una fuente real y
citable.}

\section{Análisis comparativo de soluciones existentes}
\label{sec:antecedentes-comparativa}

\pendiente{comparar las soluciones existentes en cuanto a capacidades de
análisis, grado de apertura de sus interfaces, soporte para datos de varios
atletas, opciones de personalización y posibilidades de extensión. Se recomienda
una tabla comparativa con criterios explícitos y justificados.}

\section{Literatura académica sobre métricas de rendimiento en carrera}
\label{sec:antecedentes-metricas}

\pendiente{revisar la literatura académica sobre las métricas de rendimiento en
carrera que el proyecto va a utilizar (ritmo, frecuencia cardíaca, potencia,
cadencia y cargas de entrenamiento como TRIMP, ACWR o VDOT), indicando para cada
una su definición, su fundamento fisiológico y su fuente. Ver la skill
\texttt{running-metrics} del repositorio para las definiciones canónicas.}

\section{Laguna que cubre este TFG y relación con los objetivos}
\label{sec:antecedentes-laguna}

\pendiente{formular de manera explícita la laguna que el TFG pretende cubrir a
partir de la comparativa y de la revisión anteriores, y enlazarla con los
objetivos del capítulo \ref{ch:objetivos}. Sin esta sección el capítulo queda
como una enumeración sin propósito.}
